\documentclass[conference]{IEEEtran}
\IEEEoverridecommandlockouts

\usepackage[utf8]{inputenc}
\usepackage[T1]{fontenc}
\usepackage{lmodern}
\usepackage[spanish]{babel}
\usepackage{cite}
\usepackage{amsmath,amssymb,amsfonts}
\usepackage{algorithmic}
\usepackage{graphicx}
\usepackage{textcomp}
\usepackage{xcolor}
\usepackage{booktabs}
\usepackage{url}
\usepackage{csquotes}

\def\BibTeX{{\rm B\kern-.05em{\sc i\kern-.025em b}\kern-.08em
    T\kern-.1667em\lower.7ex\hbox{E}\kern-.125emX}}

\begin{document}

\title{Resolución de Ambigüedad en Predicción Conforme\\
mediante Optimización de Explicaciones Contrafácticas\\
en Espacios de Características Tabulares}

\author{
\IEEEauthorblockN{Gonzalo Michea}
\IEEEauthorblockA{\textit{Ingeniería Civil en Informática} \\
\textit{Universidad Técnica Federico Santa María}\\
Valparaíso, Chile \\
gonzalo.michea@usm.cl}
\and
\IEEEauthorblockN{Sofía Ramírez}
\IEEEauthorblockA{\textit{Ingeniería Civil en Informática} \\
\textit{Universidad Técnica Federico Santa María}\\
Valparaíso, Chile \\
sofia.ramirez@usm.cl}
\and
\IEEEauthorblockN{Marco Repetto}
\IEEEauthorblockA{\textit{Ingeniería Civil en Informática} \\
\textit{Universidad Técnica Federico Santa María}\\
Valparaíso, Chile \\
marco.repetto@usm.cl}
}

\maketitle

\begin{abstract}
La predicción conforme (PC) ofrece un marco estadístico robusto para la
cuantificación de incertidumbre, garantizando cobertura marginal formal
bajo el único supuesto de intercambiabilidad de los datos. Sin embargo,
la aparición de conjuntos de predicción ambiguos ($|C(X)| > 1$) limita
su aplicabilidad en sistemas críticos de toma de decisiones, como el
diagnóstico clínico o la evaluación de riesgo crediticio, donde se
requiere una decisión única y justificable. Este artículo propone un
marco metodológico que integra la PC con la Inteligencia Artificial
Explicable (XAI) mediante explicaciones contrafácticas optimizadas, un
método post-hoc agnóstico al modelo que genera perturbaciones mínimas
e interpretables sobre las características de entrada. A diferencia de
enfoques anteriores, se define una función de pérdida compuesta
$\mathcal{L}(X, X')$ que incorpora explícitamente el umbral conforme
$q_\alpha$, buscando perturbaciones mínimas $X'$ tales que el puntaje
de no conformidad satisfaga estrictamente la condición de unicidad del
conjunto de predicción. Adicionalmente, se garantiza la plausibilidad
morfológica de los contrafácticos mediante restricciones de coherencia
multivariada basadas en distancia de Mahalanobis condicional por clase.
El enfoque se evalúa sobre el \textit{Dry Bean Dataset} con nivel de
confianza $1-\alpha = 0{,}95$ ($\alpha = 0{,}05$), obteniendo entre 116
y 204 instancias conformemente ambiguas según el modelo base, sobre las
cuales se aplica el proceso de resolución contrafáctica. Se hipotetiza
que la ambigüedad conforme está estrechamente relacionada con la
proximidad a las fronteras de decisión, y que el razonamiento
contrafáctico ofrece un mecanismo interpretable y accionable para
navegar dicha incertidumbre.
\end{abstract}

\begin{IEEEkeywords}
predicción conforme, explicaciones contrafácticas, cuantificación de
incertidumbre, inteligencia artificial explicable, optimización restringida,
clasificación tabular, distancia de Mahalanobis condicional.
\end{IEEEkeywords}

% ─────────────────────────────────────────────
\section{Introducción}

En numerosos dominios de alto impacto, los sistemas de clasificación
automática deben no solo predecir una etiqueta, sino también comunicar
cuándo \emph{no} están en condiciones de hacerlo con suficiente certeza.
En medicina, por ejemplo, un modelo que clasifica imágenes histológicas
puede enfrentarse a muestras que comparten rasgos de dos o más patologías;
en finanzas, un clasificador de riesgo crediticio puede encontrar
solicitantes cuyo perfil se ubica en la frontera entre distintas categorías
de riesgo. En ambos casos, una predicción puntual sin cuantificación de
incertidumbre puede inducir decisiones clínicas o económicas erróneas con
consecuencias graves. La necesidad de sistemas de predicción con garantías
estadísticas formales y capacidad de explicación es, por lo tanto, un
problema abierto y de creciente relevancia en la comunidad de aprendizaje
automático.

Los modelos de clasificación contemporáneos, a pesar de su alta precisión
empírica, carecen frecuentemente de mecanismos formales para comunicar su
propia incertidumbre. La Predicción Conforme (PC) ha surgido como solución
teórica potente~\cite{b3}, transformando predicciones puntuales en conjuntos
de predicción $C(X)$ que contienen la etiqueta verdadera con probabilidad
predefinida $1-\alpha$, bajo el único supuesto de intercambiabilidad de los
datos.

Sin embargo, en regiones del espacio de características donde las fronteras
de decisión se solapan, el modelo produce conjuntos ambiguos ($|C(X)| > 1$).
Consideremos el escenario concreto de un sistema de apoyo al diagnóstico que
devuelve $C(X) = \{\text{patología A},\, \text{patología B}\}$: el clínico
sabe que el modelo es estadísticamente correcto al 95\,\%, pero no dispone de
información accionable sobre qué característica del paciente, si cambiara,
inclinaría la predicción hacia una sola clase. Este fenómeno no debe verse
solo como un fallo del modelo, sino como una señal estructural: la instancia
$X$ se encuentra en una región de alta incertidumbre epistémica, donde
cambios pequeños e interpretables en las características pueden ser
suficientes para resolver la ambigüedad~\cite{b1}.

En este trabajo, se aborda la ambigüedad conforme como un estado de
incertidumbre \textit{resoluble} mediante perturbaciones contrafácticas~\cite{b4}.
Las explicaciones contrafácticas constituyen un método de XAI \textbf{post-hoc},
\textbf{agnóstico al modelo} y de alcance \textbf{local} (por instancia),
que responde a la pregunta: ``¿qué cambios mínimos en las características
de entrada provocarían un resultado diferente?''~\cite{b4,b7}. La
contribución principal es la formalización de un proceso de búsqueda que
busca colapsar el conjunto de predicción a un singleton $\{Y_t\}$ en el
punto contrafáctico $X'$. Si bien dicho punto pierde las garantías formales
de cobertura marginal del split conformal tradicional debido a la ruptura
de intercambiabilidad, el colapso actúa como guía heurística de certidumbre
local del modelo, orientando la búsqueda hacia regiones de baja ambigüedad
percibida para la clase objetivo.

\subsection{Contexto de los Datos}

El estudio se realiza sobre el \textbf{Dry Bean Dataset}~\cite{b5},
disponible en el repositorio UCI Machine Learning Repository, compuesto
por 13\,611 instancias de siete variedades de frijol seco descritas
mediante 16 características morfológicas continuas obtenidas mediante
visión por computadora. Los datos no requieren imputación de valores
faltantes, aunque sí un preprocesamiento de estandarización
($\mu = 0$, $\sigma = 1$) previo al entrenamiento. Las clases presentan un
leve desbalanceo (la clase \textit{Bombay} tiene 522 instancias frente
a las 3\,546 de \textit{Dermason}), lo que será considerado en el
análisis de resultados. La elección de este dataset se justifica en
detalle en la Sección~\ref{sec:datos}.

\subsection{Trabajos Relacionados}

Maalej et al.~\cite{b6} propusieron recientemente explicaciones
contrafácticas específicamente diseñadas para conjuntos de predicción
conformes, demostrando la viabilidad del enfoque en datos de texto. Este
trabajo se diferencia en tres aspectos: (1) se enfoca en datos tabulares
con características morfológicas correlacionadas; (2) incorpora
restricciones de plausibilidad multivariada mediante distancia de
Mahalanobis condicional por clase objetivo; y (3) evalúa explícitamente
la relación entre la ambigüedad y la proximidad a fronteras de decisión
mediante reducción de dimensionalidad.

Wachter et al.~\cite{b4} establecieron el marco teórico de las
explicaciones contrafácticas como mecanismo de XAI sin acceso a los
internos del modelo. En el ámbito de la cuantificación de incertidumbre,
Angelopoulos y Bates~\cite{b2} proveen una introducción exhaustiva a la
predicción conforme, mientras que Stutz et al.~\cite{b1} analizan
específicamente el problema de la ambigüedad inherente en los conjuntos
de predicción. Este trabajo conecta formalmente estas dos líneas de
investigación para datos tabulares.

\subsection{Pregunta de Investigación}

Dada una instancia $X$ tal que su conjunto de predicción bajo nivel de
confianza $1-\alpha = 0{,}95$ es ambiguo ($|C(X)| > 1$), ¿existe una
perturbación mínima $\delta$ tal que $X' = X + \delta$ sea morfológicamente
plausible respecto a la clase objetivo $Y_t$ y satisfaga $|C(X')| = 1$?

Esta pregunta motiva un estudio con tres objetivos complementarios y medibles:
\begin{itemize}
    \item[$\bullet$] \textbf{O1 --- Interpretabilidad}: determinar, mediante
    reducción de dimensionalidad (t-SNE, UMAP)~\cite{b11,b12}, que las
    instancias ambiguas identificadas bajo $\alpha = 0{,}05$ (204 para LR,
    179 para XGBoost y 116 para MLP) se concentran estadísticamente en las
    fronteras de decisión entre clases, cuantificando la proporción de
    ambigüedad por región del espacio de características.

    \item[$\bullet$] \textbf{O2 --- Accionabilidad}: lograr una tasa de éxito
    de resolución $\geq 70\%$ sobre las instancias ambiguas del
    conjunto de evaluación, definida como la proporción de instancias para
    las cuales el algoritmo encuentra un contrafáctico $X'$ con
    $|C(X')| = 1$, $d_M(X', \mathcal{D}_{Y_t}) \leq \tau_M$ y
    $\|X' - X\|_2 \leq \tau_d$.

    \item[$\bullet$] \textbf{O3 --- Comparación de modelos}: establecer si
    la complejidad del modelo base (Regresión Logística vs.\ XGBoost vs.\
    MLP) afecta significativamente la tasa de éxito y el costo de resolución,
    reportando diferencias estadísticamente significativas mediante prueba
    de Wilcoxon con $p < 0{,}05$~\cite{b13}.
\end{itemize}

La implementación completa del proyecto se encuentra disponible en:
\url{https://github.com/Sofiwiwi/Proyecto-XAI}

% ─────────────────────────────────────────────
\section{Método Propuesto}

En esta sección se describe el marco teórico, los datos utilizados, el
entorno de experimentación y la arquitectura completa de la solución que
integra Predicción Conforme con explicaciones contrafácticas para resolver
la ambigüedad en conjuntos de predicción sobre datos tabulares.\footnote{El
código y resultados se mantendrán en el repositorio público del proyecto
en GitHub, sincronizado con los notebooks ejecutados en Kaggle.}

\subsection{Marco teórico}

El modelo propuesto integra Predicción Conforme (CP) con explicaciones
contrafácticas post-hoc para resolver la ambigüedad en conjuntos de
predicción multiclase. Cuando las fronteras de decisión entre clases se
solapan, la CP genera conjuntos ambiguos con $|C(x)|>1$, lo que dificulta
la toma de decisiones en dominios críticos donde se requiere una única
etiqueta justificable~\cite{b1}. En este trabajo se interpreta dicha
ambigüedad como un estado de incertidumbre \emph{resoluble} mediante
cambios mínimos y plausibles en las características de entrada, lo que
motiva el uso de explicaciones contrafácticas como herramienta de XAI
local y accionable~\cite{b4,b7}.

Las explicaciones contrafácticas responden a la pregunta ``¿qué debería
cambiar en esta instancia para obtener una predicción diferente?'',
generando un punto $x'$ cercano a $x$ que altera la predicción del
modelo~\cite{b4}. La contribución central consiste en definir una función
de pérdida que incorpora explícitamente el umbral conforme $q_\alpha$ y
que busca contrafácticos $x'$ tales que $C(x') = \{y_t\}$, proporcionando
una etiqueta única y una explicación interpretable que sirve como guía
heurística de certidumbre local del modelo base en el punto modificado,
aun cuando las garantías estadísticas globales no se preserven en la
instancia optimizada~\cite{b1,b6}.

\subsection{Datos: Dry Bean Dataset}
\label{sec:datos}

Se utilizará el \textit{Dry Bean Dataset}~\cite{b5}, disponible en el
repositorio UCI Machine Learning Repository. Contiene 13\,611 instancias
de siete variedades de frijol seco (\textit{Seker, Barbunya, Bombay, Cali,
Dermason, Horoz} y \textit{Sira}), descritas mediante 16 características
morfológicas continuas (área, perímetro, ejes mayor y menor, entre otras)
obtenidas mediante visión por computadora. El dataset es público, no
requiere licencia especial y está libre de valores faltantes.

Este conjunto de datos es especialmente adecuado para estudiar ambigüedad
conforme por dos razones técnicas. Primero, varias clases presentan alto
solapamiento en el espacio de características; en particular,
\textit{Dermason} y \textit{Sira} comparten rangos similares en varios
ejes morfológicos, generando fronteras de decisión difusas que producen
naturalmente conjuntos de predicción con $|C(x)|>1$. Segundo, las 16
características exhiben fuertes correlaciones geométricas (por ejemplo,
área y perímetro), lo que justifica el uso de restricciones multivariadas
para garantizar plausibilidad morfológica.

El preprocesamiento consiste en una estandarización ($\mu = 0$, $\sigma = 1$)
de las características y en la partición del dataset en tres subconjuntos
disjuntos: entrenamiento (60\,\%), calibración (20\,\%) y evaluación
(20\,\%). Con $\alpha = 0{,}05$, el conjunto de evaluación genera entre
116 y 204 instancias conformemente ambiguas según el modelo base
(LR: 204, XGBoost: 179, MLP: 116), suficientes para los experimentos
de resolución contrafáctica. Dado el leve desbalanceo de clases, se
reportan métricas macro-promediadas además de las globales~\cite{b5}.

\subsection{Entorno de experimentación}

Los experimentos se ejecutaron en un notebook de Kaggle sobre un sistema
Linux \texttt{6.6.122+ x86\_64} con Python \texttt{3.12.12}, 2 núcleos
físicos, 4 lógicos y 31{,}35~GB de RAM, haciendo uso del acelerador GPU
T4x2 provisto por la plataforma. El entorno incluye librerías estándar
de aprendizaje automático preconfiguradas: \texttt{scikit-learn},
\texttt{xgboost}, \texttt{NumPy} y \texttt{SciPy}, lo que facilita la
implementación sin necesidad de infraestructura especializada.

\subsection{Arquitectura general de la solución}

El modelo propuesto se apoya en una arquitectura comparativa con tres
clasificadores base: Regresión Logística, XGBoost y una Red Neuronal
Multicapa (MLP). La Regresión Logística actúa como línea base
interpretable, mientras que XGBoost y MLP representan modelos de caja
negra con distinta capacidad de generalización, lo que permite estudiar
cómo la complejidad del modelo afecta la ambigüedad conforme y la
factibilidad de resolverla mediante contrafácticos (objetivo O3).

Para cada modelo base se aplica Predicción Conforme dividida
(\textit{split conformal prediction}). A partir del conjunto de
calibración se calcula un cuantil empírico $q_\alpha$ sobre un puntaje
de no conformidad $s(x,y)$ definido a partir de las probabilidades de
clase; este cuantil se utiliza luego para construir, en el conjunto de
evaluación, conjuntos de predicción $C(x)$ con nivel de confianza
$1-\alpha = 0{,}95$. Las instancias ambiguas se identifican como aquellas
para las que $|C(x)|>1$.

De forma resumida, el pipeline de la solución es el siguiente:
\begin{enumerate}
    \item Cargar y preprocesar el \textit{Dry Bean Dataset}
          (estandarización y particionado).
    \item Entrenar los modelos base (Regresión Logística, XGBoost, MLP).
    \item Calcular el cuantil conforme $q_\alpha$ a partir del conjunto
          de calibración para cada modelo.
    \item Construir los conjuntos de predicción $C(x)$ en el conjunto de
          evaluación e identificar las instancias ambiguas con $|C(x)|>1$.
    \item Para cada instancia ambigua, resolver un problema de
          optimización contrafáctica para encontrar un $x'$ plausible
          con $|C(x')|=1$.
    \item Evaluar la tasa de éxito, el costo de resolución y la
          plausibilidad de los contrafácticos generados.
\end{enumerate}

\subsection{Predicción Conforme dividida}

Se utiliza Predicción Conforme dividida para clasificación~\cite{b2,b3}.
Dado un conjunto de calibración $\{(x_i,y_i)\}_{i=1}^n$ y un puntaje
de no conformidad $s(x,y) = 1 - \hat{p}(y\mid x)$, se define el cuantil
empírico:
\begin{equation}
    q_{\alpha} = Q\!\left(\{s(x_i, y_i)\}_{i=1}^{n};
    \frac{\lceil (1-\alpha)(n+1) \rceil}{n}\right),
    \label{eq:qhat}
\end{equation}
donde $Q(\cdot; p)$ denota el cuantil de nivel $p$. Se utiliza
$\alpha = 0{,}05$ (cobertura marginal del 95\,\%) en todos los
experimentos; este valor fue seleccionado a partir de un análisis
exploratorio que mostró que con $\alpha = 0{,}10$ los tres modelos base
no generan instancias ambiguas en el conjunto de evaluación, dada su
alta precisión ($\approx 92\%$). Con $\alpha = 0{,}05$ se obtienen entre
116 y 204 instancias ambiguas según el modelo (LR: 204, XGBoost: 179,
MLP: 116). La cobertura marginal garantizada es del 95\,\%, lo que
representa una garantía estadística más conservadora que en
configuraciones habituales de la literatura.

El conjunto de predicción para una nueva instancia se construye como:
\begin{equation}
    C(x) = \{y : s(x, y) \leq q_{\alpha}\}.
    \label{eq:predset}
\end{equation}
Bajo el supuesto de intercambiabilidad, se garantiza la cobertura
marginal $P(Y \in C(X)) \geq 1 - \alpha$~\cite{b2}. Las instancias
ambiguas se definen como aquellas para las cuales $|C(x)| > 1$,
indicando que múltiples clases son estadísticamente consistentes con
la entrada $x$~\cite{b1}.

\subsection{Formulación de la optimización contrafáctica}

La resolución de la ambigüedad se plantea como un problema de
optimización restringida, en el que se busca una perturbación mínima
$x'$ de una instancia ambigua $x$ que cumpla simultáneamente: (i) que
el conjunto de predicción conforme colapse a un singleton con la clase
objetivo $y_t$, y (ii) que el punto contrafáctico sea morfológicamente
plausible respecto a la distribución de la clase objetivo~\cite{b4,b6}.

Para ello se define una función de pérdida compuesta:
\begin{equation}
    \begin{aligned}
        \mathcal{L}(x,x') =&\;
        \lambda \left(\|x'-x\|_1 + \|x'-x\|_2\right) \\
        &+ \mathcal{L}_{\mathrm{conf}}(x', y_t, q_\alpha) \\
        &+ \gamma \, d_M(x', \mathcal{D}_{y_t})
    \end{aligned}
    \label{eq:loss}
\end{equation}
donde:
\begin{itemize}
    \item $\|x'-x\|_1 + \|x'-x\|_2$ combina normas para promover
    escasez de cambios y suavidad simultáneamente; $\lambda>0$ pondera
    esta regularización~\cite{b4}.

    \item $\mathcal{L}_{\mathrm{conf}}$ es una penalización tipo
    \textit{hinge} sobre los puntajes de no conformidad:
    \begin{align}
        \mathcal{L}_{\mathrm{conf}}(x', y_t, q_\alpha)
        &= \max\bigl(0,\, s(x', y_t) - q_\alpha\bigr) \nonumber\\
        &\quad + \sum_{j \neq t} \max\bigl(0,\, q_\alpha - s(x', y_j)\bigr),
        \label{eq:lconf}
    \end{align}
    de modo que se fuerza la inclusión de $y_t$ en $C(x')$ y la
    exclusión del resto de las clases~\cite{b1}.

    \item $d_M(x', \mathcal{D}_{y_t})$ es la distancia de Mahalanobis
    condicional de $x'$ respecto a la distribución de la clase objetivo
    $y_t$, y $\gamma>0$ pondera la restricción de plausibilidad. Los
    hiperparámetros $\lambda$ y $\gamma$ se seleccionan mediante búsqueda
    en grilla sobre $\lambda \in \{0{,}1;\; 1{,}0;\; 10{,}0\}$ y
    $\gamma \in \{0{,}1;\; 0{,}5;\; 2{,}0\}$, evaluada sobre las
    instancias ambiguas del conjunto de calibración. Los umbrales de
    aceptación se fijan en $\tau_d = 8{,}0$ y $\tau_M = 15{,}0$; valores
    más estrictos resultaron en tasas de éxito inferiores al 30\,\% para
    modelos lineales, dado que las instancias ambiguas bajo $\alpha = 0{,}05$
    requieren perturbaciones de mayor magnitud para cruzar las fronteras
    de decisión.
\end{itemize}

\paragraph{Limitación teórica.}
Al optimizar activamente $x'$ para minimizar $\mathcal{L}(x, x')$ y
forzar $C(x') = \{y_t\}$, se rompe el supuesto de intercambiabilidad
requerido por la Predicción Conforme. Dado que $x'$ es una instancia
sintética sometida a un sesgo de selección mediante optimización, el
conjunto $C(x')$ carece de la garantía matemática de cobertura marginal
del split conformal. Por lo tanto, el colapso del conjunto conforme debe
interpretarse no como una garantía estadística sobre el contrafáctico,
sino como un \emph{criterio local de optimización} diseñado para orientar
la búsqueda hacia regiones donde el modelo base exhibe alta confianza y
baja ambigüedad respecto a la clase objetivo.

\paragraph{Estrategia de optimización.}
La pérdida~\eqref{eq:loss} se minimiza numéricamente mediante el
algoritmo L-BFGS-B implementado en \texttt{scipy.optimize.minimize} para
los tres clasificadores base. En cada instancia ambigua se optimiza
directamente sobre el vector de características $x'$, utilizando
derivadas aproximadas numéricamente a partir de las probabilidades de
clase $\hat{p}(y \mid x')$ devueltas por cada modelo. Para garantizar
la viabilidad computacional, cada llamada al optimizador se limita a un
máximo de 200 iteraciones (\texttt{maxiter = 200}), una tolerancia de
función (\texttt{ftol = 1e-7}) y un \textit{timeout} de 10 segundos por
instancia, detenido mediante una envoltura temporal externa. Estos valores
buscan equilibrar costo computacional y calidad de los candidatos: el
criterio de aceptación es discreto ($C(x') = \{y_t\}$ o no lo es) y, por
tanto, no se requiere convergencia de alta precisión.

\subsection{Plausibilidad y coherencia multivariada}

Para garantizar que los contrafácticos generados sean morfológicamente
posibles no es suficiente restringir cada característica a su rango
univariado observado. Las 16 características del dataset exhiben
fuertes correlaciones estructurales (por ejemplo, el área crece
aproximadamente de forma cuadrática con los ejes principales), por lo
que una perturbación que viole estas relaciones produciría instancias
físicamente incoherentes.

Se incorpora la distancia de Mahalanobis \emph{condicional}
$d_M(x', \mathcal{D}_{y_t})$ como restricción suave en la función de
pérdida~\eqref{eq:loss}, donde $\mathcal{D}_{y_t}$ representa la
distribución empírica de los datos de entrenamiento correspondientes
a la clase objetivo $y_t$, caracterizada por su media $\mu_{y_t}$ y
su matriz de covarianza $\Sigma_{y_t}$:
\begin{equation}
    d_M(x', \mathcal{D}_{y_t}) =
    \sqrt{(x' - \mu_{y_t})^\top \Sigma_{y_t}^{-1} (x' - \mu_{y_t})}.
    \label{eq:maha}
\end{equation}

El uso de la métrica condicionada por clase evita penalizar
contrafácticos plausibles que se alejan de la media global debido a la
heterogeneidad física entre clases (por ejemplo, entre \textit{Bombay}
y \textit{Dermason}, que difieren considerablemente en tamaño). Un valor
bajo de $d_M$ indica que $x'$ es coherente con la morfología y
variabilidad multivariada de la clase objetivo $y_t$. En los experimentos
se analizará la distribución de esta distancia condicional antes y después
de la generación de contrafácticos para verificar que la plausibilidad se
mantiene dentro de los rangos característicos de la clase objetivo~\cite{b10}.

% ─────────────────────────────────────────────
\section{Plan de Evaluación}

\subsection{Métricas de Rendimiento}

El enfoque propuesto será evaluado mediante las siguientes métricas:
\begin{itemize}
    \item[$\bullet$] \textbf{Tasa de éxito (Resolución Conforme)}:
    proporción de instancias ambiguas para las cuales se encontró un
    contrafáctico $X'$ con $|C(X')| = 1$ (umbral objetivo: $\geq 70\%$).

    \item[$\bullet$] \textbf{Costo de resolución}: magnitud de la
    perturbación necesaria, medida en normas $L_1$ y $L_2$.

    \item[$\bullet$] \textbf{Preservación de plausibilidad}: distancia
    de Mahalanobis condicional promedio de los contrafácticos respecto
    a la distribución de la clase objetivo $\mathcal{D}_{y_t}$.

    \item[$\bullet$] \textbf{Reducción del conjunto de predicción}:
    cambio en $|C(X)|$ antes y después de la perturbación.

    \item[$\bullet$] \textbf{Análisis de fronteras de decisión}:
    visualización de las regiones de ambigüedad conforme mediante
    reducción de dimensionalidad con t-SNE y UMAP~\cite{b11,b12}, para
    verificar que los contrafácticos cruzan hacia regiones no ambiguas.
\end{itemize}

\subsection{Evaluación de Explicaciones Contrafácticas}

Para evaluar la calidad de las explicaciones generadas se emplea un
conjunto de métricas nativas y teóricamente consistentes con métodos
basados en ejemplos~\cite{b8,b9,b10}, en lugar de métricas de fidelidad
de atribución que asumen vectores de importancia y resultan
conceptualmente incompatibles con contrafácticos:

\begin{itemize}
    \item[$\bullet$] \textbf{Escasez (Sparsity)}~\cite{b8}: número de
    características modificadas respecto a la instancia original,
    $S(x, x') = \sum_{j=1}^d \mathbb{I}(|x'_j - x_j| > \epsilon)$,
    donde $\epsilon$ es un umbral de tolerancia numérica. Se persigue
    alta escasez para maximizar la interpretabilidad de la explicación.

    \item[$\bullet$] \textbf{Proximidad (Proximity)}~\cite{b4,b9}:
    distancias $L_1$ y $L_2$ normalizadas por la desviación estándar de
    cada variable en el conjunto de entrenamiento, cuantificando cuánto
    cambia el contrafáctico respecto al original.

    \item[$\bullet$] \textbf{Plausibilidad (Plausibility)}~\cite{b10}:
    evaluada cuantitativamente a través de $d_M(x', \mathcal{D}_{y_t})$
    (ecuación~\eqref{eq:maha}), garantizando que el punto generado sea
    físicamente consistente con la clase objetivo.

    \item[$\bullet$] \textbf{Accionabilidad (Actionability)}: se
    identificará \emph{a priori} el subconjunto de características
    morfológicamente accionables (aquellas que pueden variar
    independientemente, como elongación y compacidad) en contraposición
    a las derivadas (como área o perímetro). Se reportará la proporción
    de contrafácticos que modifican exclusivamente características
    accionables como indicador de la utilidad práctica de la explicación.
\end{itemize}

\subsection{Comparación entre Modelos}

Se comparará el rendimiento del marco propuesto entre los tres
clasificadores base (XGBoost, MLP y Regresión Logística), reportando
para cada uno la tasa de éxito, el costo de resolución promedio y la
plausibilidad de los contrafácticos generados. Las diferencias se
evaluarán con prueba de Wilcoxon con nivel de significancia
$p < 0{,}05$~\cite{b13}.

\subsection{Factibilidad Computacional}

El proyecto es factible dentro del semestre: el Dry Bean Dataset es de
acceso libre~\cite{b5}; la predicción conforme dividida se implementa de
forma nativa sin dependencias adicionales; y la optimización contrafáctica
se realiza mediante \texttt{scipy.optimize.minimize} con L-BFGS-B,
utilizando un máximo de 200 iteraciones y un \textit{timeout} de
10 segundos por instancia. Las métricas de evaluación se calculan con
\texttt{NumPy} y \texttt{SciPy}. Todo el pipeline puede ejecutarse en
CPU estándar sin infraestructura especializada, como la proporcionada
por Kaggle.

% ─────────────────────────────────────────────
\section{Resultados Experimentales}

En esta sección se reportan los resultados preliminares obtenidos tras
la primera ejecución completa del pipeline experimental. Los resultados
presentados corresponden a una validación temprana del flujo de
entrenamiento, calibración conforme, detección de ambigüedad y resolución
contrafáctica, y deben interpretarse como preliminares.

\subsection{Entorno de ejecución}

Los experimentos fueron ejecutados en Kaggle sobre un sistema Linux
\texttt{6.6.122+ x86\_64} con Python \texttt{3.12.12} (GCC
\texttt{11.4.0}), 2 núcleos físicos, 4 lógicos, 31{,}35~GB de RAM y
acelerador GPU T4x2. Se utilizaron \texttt{NumPy}, \texttt{Pandas},
\texttt{Matplotlib}, \texttt{SciPy}, \texttt{Scikit-learn},
\texttt{XGBoost} y \texttt{ucimlrepo}.

\subsection{Resultados de clasificación}

La Tabla~\ref{tab:clasificacion} resume el desempeño de los tres modelos
base sobre el conjunto de prueba.

\begin{table}[h]
\centering
\caption{Resultados de clasificación sobre el conjunto de prueba.}
\label{tab:clasificacion}
\begin{tabular}{lccc}
\toprule
\textbf{Modelo} & \textbf{Accuracy} & \textbf{Macro F1} & \textbf{Weighted F1} \\
\midrule
Regresión Logística & 0{,}9199 & 0{,}93 & 0{,}92 \\
XGBoost             & 0{,}9203 & 0{,}93 & 0{,}92 \\
MLP                 & 0{,}9236 & 0{,}94 & 0{,}92 \\
\bottomrule
\end{tabular}
\end{table}

Los tres modelos obtuvieron rendimiento similar ($\approx 92\%$). MLP
presentó el mejor desempeño global (0{,}9236), seguido por XGBoost
(0{,}9203) y Regresión Logística (0{,}9199).

\subsection{Predicción conforme y detección de ambigüedad}

Se aplicó predicción conforme dividida con $\alpha = 0{,}05$. Los
umbrales conformes obtenidos fueron $\hat{q} = 0{,}6799$ (LR),
$\hat{q} = 0{,}7076$ (XGBoost) y $\hat{q} = 0{,}6769$ (MLP).
La Tabla~\ref{tab:conforme} resume los resultados.

\begin{table}[t]
\centering
\caption{Resultados de predicción conforme ($\alpha = 0{,}05$).}
\label{tab:conforme}
\resizebox{\columnwidth}{!}{%
\begin{tabular}{lccccc}
\toprule
\textbf{Modelo} & \textbf{Cobertura} & \textbf{N ambig.} &
\textbf{\% ambig.} & $\mathbf{|C|}$ \textbf{prom.} & \textbf{Vacíos} \\
\midrule
Regresión Logística & 0{,}9479 & 204 & 7{,}5\% & 1{,}075 & 0 \\
XGBoost             & 0{,}9486 & 179 & 6{,}6\% & 1{,}066 & 0 \\
MLP                 & 0{,}9412 & 116 & 4{,}3\% & 1{,}043 & 0 \\
\bottomrule
\end{tabular}}
\end{table}

Las coberturas empíricas se mantuvieron cercanas al 95\,\% en los tres
modelos. Ningún modelo generó conjuntos vacíos. La menor tasa de
ambigüedad en MLP (4{,}3\,\%) refleja que sus fronteras de decisión más
complejas producen predicciones de mayor certeza bajo esta cobertura, en
contraste con las fronteras lineales de LR (7{,}5\,\%).

\subsection{Análisis de ambigüedad por par de clases}

La Tabla~\ref{tab:pares} muestra los pares de clases más frecuentes en
los conjuntos ambiguos de LR. El par \textit{DERMASON}--\textit{SIRA}
concentra 113 de las 204 instancias ambiguas (55{,}4\,\%), confirmando el
alto solapamiento morfológico entre estas variedades descrito en la
Sección~\ref{sec:datos}.

\begin{table}[h]
\centering
\caption{Pares de clases más frecuentes en conjuntos ambiguos (LR).}
\label{tab:pares}
\begin{tabular}{lc}
\toprule
\textbf{Par de clases} & \textbf{N instancias} \\
\midrule
DERMASON -- SIRA  & 113 \\
SEKER -- SIRA     & 23  \\
BARBUNYA -- CALI  & 23  \\
HOROZ -- SIRA     & 15  \\
CALI -- HOROZ     & 8   \\
CALI -- SIRA      & 7   \\
DERMASON -- SEKER & 6   \\
BARBUNYA -- SIRA  & 5   \\
\bottomrule
\end{tabular}
\end{table}

\subsection{Resolución contrafáctica}

La Tabla~\ref{tab:contrafacticos} resume los resultados de la primera
ejecución completa del optimizador contrafáctico.

\begin{table}[t]
\centering
\caption{Resultados de resolución contrafáctica.}
\label{tab:contrafacticos}
\resizebox{\columnwidth}{!}{%
\begin{tabular}{lcccccc}
\toprule
\textbf{Modelo} & \textbf{N ambig.} & \textbf{Res.} & \textbf{Éxito} &
\textbf{L2 prom.} & $\mathbf{d_M}$ \textbf{prom.} & \textbf{T total} \\
\midrule
LR      & 204 & 182 & 89{,}2\% & 0{,}3817 & 3{,}0098 & 34{,}0 min \\
XGBoost & 179 & 3   & 1{,}7\%  & 0{,}2389 & 2{,}9971 & 29{,}9 min \\
MLP     & 116 & 114 & 98{,}3\% & 0{,}2452 & 2{,}7679 & 19{,}3 min \\
\bottomrule
\end{tabular}}
\end{table}

LR y MLP superaron ampliamente el umbral del 70\,\% definido en O2, con
tasas de éxito de 89{,}2\,\% y 98{,}3\,\% respectivamente. XGBoost
resolvió únicamente 3 de 179 instancias (1{,}7\,\%), lo que evidencia
una limitación estructural de L-BFGS-B aplicado a modelos basados en
árboles: sus fronteras de decisión constantes por trozos producen
gradientes numéricos nulos en la mayoría del dominio, impidiendo que el
optimizador encuentre direcciones de descenso efectivas. Este resultado
contrasta marcadamente con el desempeño en LR y MLP, cuyos modelos
suaves permiten la propagación de gradientes. La diferencia entre modelos
constituye evidencia directa para el Objetivo O3 y será analizada
estadísticamente mediante prueba de Wilcoxon en la evaluación final.

La distancia de Mahalanobis condicional promedio de los contrafácticos
($d_M \approx 3{,}0$) es comparable con la media de referencia calculada
sobre los datos reales de entrenamiento (3{,}3419), lo que indica que
los contrafácticos generados son morfológicamente plausibles respecto a
la distribución de la clase objetivo.

\subsection{Observaciones preliminares}

Los resultados iniciales validan el funcionamiento general del pipeline.
Los modelos base alcanzan desempeños similares, la predicción conforme
produce coberturas cercanas al nivel objetivo, y LR y MLP demuestran
que la resolución de ambigüedad es factible con el optimizador propuesto.
La baja tasa de XGBoost anticipa la necesidad de una estrategia de
optimización alternativa para modelos no diferenciables, aspecto que
será abordado en la evaluación final.

% ─────────────────────────────────────────────
\begin{thebibliography}{00}

\bibitem{b1} D. Stutz, A. G. Roy, T. Matejovicova, P. Strachan,
A. T. Cemgil, y A. Doucet, ``Conformal prediction under ambiguous
ground truth,'' \textit{arXiv preprint} arXiv:2307.09302, 2023.

\bibitem{b2} A. N. Angelopoulos y S. Bates, ``A gentle introduction to
conformal prediction and distribution-free uncertainty quantification,''
\textit{arXiv preprint} arXiv:2107.07511, 2022.

\bibitem{b3} R. Tibshirani, ``Conformal prediction,'' Notas de clase,
University of California Berkeley, 2023.

\bibitem{b4} S. Wachter, B. Mittelstadt, y C. Russell,
``Counterfactual explanations without opening the black box,''
\textit{Harvard Journal of Law \& Technology}, vol.~31, no.~2, 2018.

\bibitem{b5} M. Koklu e I. A. Ozkan, ``Multiclass classification of
dry beans using computer vision and machine learning techniques,''
\textit{Computers and Electronics in Agriculture}, vol.~174,
p.~105507, 2020.

\bibitem{b6} R. Maalej, R. Confalonieri, y M. Lippi,
``Counterfactual explanations for conformal prediction sets,''
in \textit{Proc. PMLR}, vol.~266, 2025.

\bibitem{b7} C. Molnar, \textit{Interpretable Machine Learning: A Guide
for Making Black Box Models Explainable}, 2nd ed., 2022. [En línea].
Disponible: \url{https://christophm.github.io/interpretable-ml-book/}

\bibitem{b8} R. Guidotti, ``Counterfactual explanations and adversarial
examples: Common grounds, mysteries, and perspectives,''
\textit{Frontiers in Computer Science}, vol.~3, p.~611244, 2021.

\bibitem{b9} A.-H. Karimi, G. Barthe, B. Schölkopf, e I. Valera,
``A survey of algorithmic recourse: definitions, formulations, state of
the art, and future directions,'' \textit{ACM SIGKDD Explorations
Newsletter}, vol.~24, no.~1, pp.~46--48, 2022.

\bibitem{b10} M. T. Keane y B. Smyth, ``Good counterfactuals and where
to find them: A case-based technique for generating explanations for
black-box models,'' in \textit{Proc. Case-Based Reasoning Research and
Development}, pp.~163--178, 2020.

\bibitem{b11} L. van der Maaten y G. Hinton, ``Visualizing data using
t-SNE,'' \textit{Journal of Machine Learning Research}, vol.~9,
pp.~2579--2605, 2008.

\bibitem{b12} L. McInnes, J. Healy, y J. Melville, ``UMAP: Uniform
Manifold Approximation and Projection for Dimension Reduction,''
\textit{arXiv preprint} arXiv:1802.03426, 2018.

\bibitem{b13} F. Wilcoxon, ``Individual comparisons by ranking methods,''
\textit{Biometrics Bulletin}, vol.~1, no.~6, pp.~80--83, 1945.

\end{thebibliography}

\end{document}